\documentclass[a4paper,11pt]{article}
\usepackage{exam-style-341}
\examinehead{341 农业综合知识三 · 网络技术分章练习}{第06章 网络层}
\begin{document}

\unitcover{341 农业综合知识三 · 网络技术分章练习}{第06章\quad 网络层}{练习（共 100 分）}{约 75 分钟}{IP 数据报格式、IP 地址分类与私有地址、划分子网与子网掩码、CIDR 与路由聚合、IP 分片、ARP／ICMP／IGMP、RIP／OSPF／BGP、最长前缀匹配、IPv6 与 NAT。}

\begin{knowmap}
\textbf{网络层服务：}无连接、尽最大努力交付；IP 数据报＝首部（固定 20 字节）＋数据，不保证可靠、不保证按序\\[2pt]
\textbf{地址体系：}分类地址 A／B／C／D／E $\rightarrow$ 划分子网（子网掩码）$\rightarrow$ CIDR（网络前缀／斜线记法）$\rightarrow$ 路由聚合\\[2pt]
\textbf{关键计算：}可用子网数 $=2^{n}-2$（$n$ 为借位数，经典方法）；可用主机数 $=2^{\text{主机位数}}-2$；片偏移＝该片数据起点$\,/\,8$；每片最大数据＝MTU$-$20\\[2pt]
\textbf{协议族：}ARP（IP$\rightarrow$MAC）· ICMP（差错报告与询问）· IGMP（IP 多播）· NAT（私有地址$\leftrightarrow$公有地址）\\[2pt]
\textbf{路由协议：}RIP（距离向量，跳数 $\le 15$）· OSPF（链路状态，Dijkstra 算法）· BGP（路径向量，自治系统之间）\\[2pt]
\textbf{地址演进：}IPv4（32 位）$\rightarrow$ IPv6（128 位，简化首部、支持自动配置）
\end{knowmap}

\begin{keyterms}
\textbf{IP 数据报}：网络层传输的数据单元，由固定 20 字节首部、可选字段和数据部分组成。\\
\textbf{MTU}：最大传送单元，数据链路层帧中数据字段的最大长度，超过时 IP 层必须分片。\\
\textbf{片偏移}：分片后某片数据在原数据报数据部分中的相对位置，以 8 字节为单位。\\
\textbf{子网掩码}：32 位二进制数，网络号与子网号对应位为 1、主机号对应位为 0。\\
\textbf{CIDR}：无分类域间路由选择，用“网络前缀／前缀长度”表示地址块，不再区分 A、B、C 类。\\
\textbf{路由聚合}：把若干个连续的小 CIDR 地址块合并为一个较大的地址块，以减少路由表项。\\
\textbf{最长前缀匹配}：查找转发表时，若多个表项都匹配，选择前缀最长（最具体）的那一项。\\
\textbf{ARP}：地址解析协议，在同一局域网内把 IP 地址解析为 MAC 地址。\\
\textbf{ICMP}：网际控制报文协议，用于报告差错和提供询问，报文封装在 IP 数据报中。\\
\textbf{RIP}：内部网关协议，基于距离向量算法，以跳数为度量，最大跳数 15。\\
\textbf{OSPF}：内部网关协议，基于链路状态算法，用 Dijkstra 算法计算最短路径。\\
\textbf{BGP}：外部网关协议，基于路径向量算法，用于自治系统之间的路由选择。\\
\textbf{NAT}：网络地址转换，在私有地址与公有地址之间进行转换，缓解 IPv4 地址短缺。\\
\textbf{IPv6}：下一代网际协议，地址长度 128 位，首部简化，支持自动配置。
\end{keyterms}

\begin{contrast}{分类地址 / 划分子网 / CIDR（每行依次为：地址结构·掩码或前缀表示·主要优点）}
\citem{分类地址}{网络号＋主机号两级结构，按首字节分为 A、B、C 类}{使用默认掩码，如 C 类为 255.255.255.0}{编址简单，但地址利用率低}
\citem{划分子网}{网络号＋子网号＋主机号三级结构}{从主机号借位作子网号，掩码中 1 的个数增加}{可在单位内部灵活划分，减少地址浪费}
\citem{CIDR}{网络前缀＋主机号，前缀长度任意}{斜线记法，如 202.117.35.0/27}{可按需分配、支持路由聚合，提高利用率}
\end{contrast}

\begin{contrast}{RIP / OSPF / BGP（每行依次为：协议类型·算法与度量·适用范围）}
\citem{RIP}{内部网关协议（IGP）}{距离向量算法，以跳数为度量，最大 15 跳}{小型网络，收敛慢}
\citem{OSPF}{内部网关协议（IGP）}{链路状态算法，Dijkstra 计算最短路径}{大型网络，收敛快、支持分层}
\citem{BGP}{外部网关协议（EGP）}{路径向量算法，在自治系统之间交换可达性信息}{互联网各自治系统之间}
\end{contrast}

\begin{contrast}{ARP / ICMP / IGMP（每行依次为：主要作用·报文封装·典型用途）}
\citem{ARP}{把 IP 地址解析为 MAC 地址}{封装在以太网帧中直接传送}{同一局域网内主机或路由器的地址解析}
\citem{ICMP}{报告差错、提供询问}{封装在 IP 数据报中}{目的不可达、超时、回送请求与回答（ping）}
\citem{IGMP}{管理 IP 多播组的成员}{封装在 IP 数据报中}{主机加入或退出多播组}
\end{contrast}

\qtypename{一、单项选择题}{10 题，每题 3 分，共 30 分}
\q IPv4 数据报首部的固定部分长度是（\quad）
\choicesline{A. 8 字节}{B. 20 字节}{C. 40 字节}{D. 60 字节}

\q IP 数据报首部中“片偏移”字段的计量单位是（\quad）
\choicesline{A. 1 字节}{B. 4 字节}{C. 8 字节}{D. 16 字节}

\q 下列 IP 地址中属于私有地址的是（\quad）
\choicesline{A. 172.15.5.1}{B. 172.16.5.1}{C. 172.32.5.1}{D. 192.169.5.1}

\q C 类 IP 地址的默认子网掩码是（\quad）
\choicesline{A. 255.0.0.0}{B. 255.255.0.0}{C. 255.255.255.0}{D. 255.255.255.255}

\q 在同一局域网内，把 IP 地址解析为 MAC 地址的协议是（\quad）
\choicesline{A. ARP}{B. ICMP}{C. IGMP}{D. BGP}

\q ICMP 报文封装在（\quad）中传送。
\choicesline{A. IP 数据报}{B. UDP 用户数据报}{C. TCP 报文段}{D. 以太网帧}

\q RIP 协议允许的最大跳数是（\quad），跳数为 16 时表示不可达。
\choicesline{A. 8}{B. 15}{C. 16}{D. 32}

\q OSPF 协议采用的路由算法是（\quad）
\choices{A. 距离向量算法}{B. 链路状态算法}{C. 路径向量算法}{D. 随机路由算法}

\q CIDR 记法 192.168.1.0/26 中，“/26”表示（\quad）
\choicesline{A. 主机号为 26 位}{B. 网络前缀为 26 位}{C. 子网数为 26}{D. 地址总数为 26}

\q IPv6 地址的长度是（\quad）
\choicesline{A. 32 位}{B. 64 位}{C. 128 位}{D. 256 位}

\qtypename{二、多项选择题}{3 题，每题 4 分，共 12 分；每题有两个或两个以上正确项}
\q 下列属于 IP 数据报首部固定部分字段的有（\quad）
\choicesline{A. 版本}{B. 总长度}{C. 生存时间}{D. 目的端口}

\q 下列属于私有地址的有（\quad）
\choicesline{A. 10.1.2.3}{B. 172.16.0.1}{C. 192.168.1.1}{D. 172.32.0.1}

\q 下列属于路由选择协议的有（\quad）
\choicesline{A. RIP}{B. OSPF}{C. BGP}{D. ARP}

\qtypename{三、判断题}{6 题，每题 2 分，共 12 分；正确的画“√”，错误的画“×”并写出改正后的表述}
\q IP 协议提供面向连接的可靠传输服务。\quad（\quad）\quad 改正：\fixblank

\q 划分子网后，子网掩码中 1 的个数比该类的默认掩码多。\quad（\quad）\quad 改正：\fixblank

\q CIDR 路由聚合要求被聚合的地址块连续，且起始地址按聚合块的大小对齐。\quad（\quad）\quad 改正：\fixblank

\q 片偏移字段以字节为单位。\quad（\quad）\quad 改正：\fixblank

\q RIP 收敛速度快，适合用于大型网络。\quad（\quad）\quad 改正：\fixblank

\q 默认路由 0.0.0.0/0 的网络前缀长度为 0。\quad（\quad）\quad 改正：\fixblank

\qtypename{四、计算题}{1 题，共 6 分；要求写出二进制或掩码推导过程}
\q 某单位分到一个 C 类网络地址 202.117.35.0/24，现需划分为 5 个可用子网（按经典划分子网方法，子网号全 0 和全 1 不使用）。请回答：
\begin{enumerate}[leftmargin=2em, itemsep=1pt, topsep=2pt]
  \item 需从主机号中借几位作为子网号？可用子网数是多少？
  \item 子网掩码是多少？写出其二进制形式。
  \item 每个子网最多可容纳多少台主机？
  \item 第 2 个可用子网的网络地址、广播地址和可用主机地址范围各是什么？
\end{enumerate}
\anslines{5cm}

\newpage
\qtypename{五、简答题}{2 题，每题 10 分，共 20 分}
\q 简述 IP 数据报首部固定部分的主要字段及其作用（至少写出 8 个字段）。\anslines{4cm}

\q 简述 RIP、OSPF、BGP 三种路由选择协议的基本思想，并比较它们的适用场合。\anslines{4cm}

\qtypename{六、综合题}{1 题，共 20 分}
\q 某单位从 ISP 申请到地址块 211.68.4.0/22（即 4 个连续的 C 类地址块 211.68.4.0/24—211.68.7.0/24），并计划在此基础上建设内部网络。请回答：
\begin{enumerate}[leftmargin=2em, itemsep=1pt, topsep=2pt]
  \item 写出该地址块的网络前缀长度、子网掩码（点分十进制）、包含的 IP 地址总数和地址范围（6 分）。
  \item 现要把其中 211.68.4.0/24 划分为 4 个等长子网供 4 个部门使用（允许使用全 0 和全 1 子网号），写出借位数、子网掩码、每个子网可用主机数，以及第 1 个子网的网络地址和可用地址范围（8 分）。
  \item 一个总长度为 3020 字节的 IP 数据报（首部 20 字节、数据 3000 字节）经 MTU 为 1500 字节的链路转发（DF=0），求分片数、各片的数据长度、片偏移、MF 标志和总长度（6 分）。
\end{enumerate}
\anslines{8cm}

\end{document}
